% Options for packages loaded elsewhere
\PassOptionsToPackage{unicode}{hyperref}
\PassOptionsToPackage{hyphens}{url}
\PassOptionsToPackage{dvipsnames,svgnames,x11names}{xcolor}
%
\documentclass[
  11pt,
  a4paper,
]{ltjsarticle}
\usepackage{amsmath,amssymb}
\usepackage{iftex}
\ifPDFTeX
  \usepackage[T1]{fontenc}
  \usepackage[utf8]{inputenc}
  \usepackage{textcomp} % provide euro and other symbols
\else % if luatex or xetex
  \usepackage{unicode-math} % this also loads fontspec
  \defaultfontfeatures{Scale=MatchLowercase}
  \defaultfontfeatures[\rmfamily]{Ligatures=TeX,Scale=1}
\fi
\usepackage{lmodern}
\ifPDFTeX\else
  % xetex/luatex font selection
\fi
% Use upquote if available, for straight quotes in verbatim environments
\IfFileExists{upquote.sty}{\usepackage{upquote}}{}
\IfFileExists{microtype.sty}{% use microtype if available
  \usepackage[]{microtype}
  \UseMicrotypeSet[protrusion]{basicmath} % disable protrusion for tt fonts
}{}
\makeatletter
\@ifundefined{KOMAClassName}{% if non-KOMA class
  \IfFileExists{parskip.sty}{%
    \usepackage{parskip}
  }{% else
    \setlength{\parindent}{0pt}
    \setlength{\parskip}{6pt plus 2pt minus 1pt}}
}{% if KOMA class
  \KOMAoptions{parskip=half}}
\makeatother
\usepackage{xcolor}
\usepackage[margin=24mm]{geometry}
\setlength{\emergencystretch}{3em} % prevent overfull lines
\providecommand{\tightlist}{%
  \setlength{\itemsep}{0pt}\setlength{\parskip}{0pt}}
\setcounter{secnumdepth}{5}
\ifLuaTeX
  \usepackage{selnolig}  % disable illegal ligatures
\fi
\IfFileExists{bookmark.sty}{\usepackage{bookmark}}{\usepackage{hyperref}}
\IfFileExists{xurl.sty}{\usepackage{xurl}}{} % add URL line breaks if available
\urlstyle{same}
\hypersetup{
  colorlinks=true,
  linkcolor={blue},
  filecolor={Maroon},
  citecolor={Blue},
  urlcolor={blue},
  pdfcreator={LaTeX via pandoc}}

\author{}
\date{}

\begin{document}

\hypertarget{ux89e3ux7b54-05-ux884cux5217ux5f0f}{%
\section{解答 05 --- 行列式}\label{ux89e3ux7b54-05-ux884cux5217ux5f0f}}

\hypertarget{section}{%
\subsection{1}\label{section}}

\texttt{det(A)=6} なので面積は6倍です。

\hypertarget{section-1}{%
\subsection{2}\label{section-1}}

\texttt{det(A)=0} は線形変換が n
次元体積を0へ潰すことを意味します。したがって像の次元は n 未満で、列は n
本の独立方向を作れません。

\hypertarget{section-2}{%
\subsection{3}\label{section-2}}

\texttt{B} が体積を \texttt{det(B)} 倍し、その後 \texttt{A} が
\texttt{det(A)}
倍するので、合成は積になります。符号も向きの反転回数として積で合成されます。

\hypertarget{section-3}{%
\subsection{4}\label{section-3}}

\[
\det A=\det U\det\Sigma\det V^T.
\]

直交行列の行列式は ±1 なので絶対値は1です。\texttt{det\ Σ=∏σ\_i}
だから結論が得られます。

\hypertarget{section-4}{%
\subsection{5}\label{section-4}}

同じではありません。例えば \texttt{A=10\^{}\{-6\}I}
は行列式が非常に小さい一方、条件数は1で完全に well-conditioned
です。悪条件性は絶対的な体積ではなく、方向ごとの伸縮率の比
\texttt{σ\_max/σ\_min} に関係します。

\end{document}
